\documentclass[11pt,a4paper]{article}

% ==================== PACKAGES ====================
\usepackage[margin=1in]{geometry}
\usepackage{setspace}
\usepackage{titlesec}
\usepackage{enumitem}
\usepackage{hyperref}
\usepackage{xcolor}
\usepackage{fancyhdr}
\usepackage{tabularx}
\usepackage{booktabs}
\usepackage{parskip}

% ==================== COLOR SCHEME ====================
\definecolor{ThemeBlue}{RGB}{0, 62, 116}
\definecolor{DarkGray}{RGB}{51, 51, 51}
\definecolor{LightGray}{RGB}{245, 245, 245}
\definecolor{AccentOrange}{RGB}{214, 96, 77}

% ==================== HYPERLINK SETUP ====================
\hypersetup{
    colorlinks=true,
    linkcolor=ThemeBlue,
    urlcolor=ThemeBlue,
    citecolor=ThemeBlue,
    pdfborder={0 0 0}
}

% ==================== SECTION FORMATTING ====================
\titleformat{\section}
    {\Large\bfseries\color{ThemeBlue}}
    {\thesection}{1em}{}
    [\titlerule]

\titleformat{\subsection}
    {\large\bfseries\color{ThemeBlue}}
    {\thesubsection}{1em}{}

\titleformat{\subsubsection}
    {\normalsize\bfseries}
    {\thesubsubsection}{1em}{}

% ==================== HEADER AND FOOTER ====================
\pagestyle{fancy}
\fancyhf{}
\fancyhead[L]{\small\textit{Big Data in Finance}}
\fancyhead[R]{\small\textit{MSc Finance}}
\fancyfoot[C]{\thepage}
\renewcommand{\headrulewidth}{0.4pt}

% ==================== CUSTOM COMMANDS ====================
\newcommand{\coursetitle}[1]{\textbf{\LARGE\color{ThemeBlue}#1}}
\newcommand{\highlight}[1]{\textcolor{AccentOrange}{\textbf{#1}}}

% Reading list formatting
\newcommand{\recommended}[1]{\item[$\bullet$] #1}
\newcommand{\advanced}[1]{\item[$\circ$] #1}

% Week box - simple colored section header
\newcommand{\weekheader}[1]{%
    \vspace{1em}
    \noindent{\large\bfseries\color{ThemeBlue} #1}
    \vspace{0.3em}
    \hrule height 0.5pt
    \vspace{0.5em}
}

% ==================== BEGIN DOCUMENT ====================
\begin{document}

% ==================== TITLE ====================
\begin{center}
    \coursetitle{Big Data in Finance}\\[0.5em]
    \normalsize Academic Year 2025--2026\\[1em]
    
    \begin{tabular}{ll}
        \toprule
        \textbf{Programme:} & MSc Finance (Elective) \\
        \textbf{Duration:} & 7 weeks \\
        \textbf{Format:} & 3-hour lectures + 1-hour tutorials \\
        \textbf{Assessment:} & 10\% Participation, 25\% Group Project, 65\% Exam \\
        \bottomrule
    \end{tabular}
\end{center}

\vspace{1em}

% ==================== COURSE TEAM ====================
\section*{Course Team}

\begin{tabular}{@{}ll}
\textbf{Module Lead:} & Dr.\ Daniele Bianchi \\
& Email: \href{mailto:d.bianchi@qmul.ac.uk}{d.bianchi@qmul.ac.uk} \\
& Website: \url{whitesphd.com} \\
& Office hours: by appointment \\
\end{tabular}

% ==================== COURSE OVERVIEW ====================
\section{Course Overview}

\textbf{Core objective:} Apply state-of-the-art machine learning methods to financial decision-making.

This module provides a comprehensive and rigorous introduction to machine learning methods for finance, with a systematic focus on both theoretical foundations and practical implementation. The module is structured around two main applications:

\begin{itemize}[leftmargin=*]
    \item \highlight{Market timing:} Time-series prediction of aggregate market returns (with factor investing as a related cross-sectional application)
    \item \highlight{Credit risk assessment:} Default prediction and credit scoring
\end{itemize}

\noindent\textbf{Key features:} Theory + Practice, Python implementation, production deployment considerations, AI-assisted learning, real financial datasets.

\subsection{Lecture Structure}

Each \textbf{3-hour week} is divided into two parts:

\begin{itemize}[leftmargin=*]
    \item \textbf{First 1.5 hours:} Theoretical foundations, methodological insights, algorithmic procedures, mathematical concepts
    \item \textbf{Second 1.5 hours:} Financial applications, case studies, implementation details, best practices
\end{itemize}

This bipartite structure ensures students develop both conceptual understanding and practical competence with industry relevance.

\subsection{Tutorial Structure}

\textbf{1-hour hands-on sessions} focusing on implementation: guided coding exercises, production-ready code practices, clean/documented/version-controlled code, and model deployment considerations.

\textbf{AI-Assisted Learning:} Students are \emph{encouraged} to leverage AI tools (ChatGPT, Claude, etc.) for debugging code, code optimization, interpreting model outputs, and understanding error messages. \textbf{Important:} You must understand and explain all code and implementations, regardless of source.

\subsection{Learning Outcomes}

By the end of this module, you will be able to:

\begin{enumerate}[leftmargin=*]
    \item Describe complex models relevant to financial decision making
    \item Compute complex models to predict default in corporate and retail credit markets
    \item Appraise a range of quantitative techniques to build trading strategies based on stock characteristics
    \item Produce deployment-ready pipelines for investment and credit risk management
\end{enumerate}

\subsection{Prerequisites}

\begin{itemize}[leftmargin=*]
    \item Basic Python programming
    \item Statistics and probability theory
    \item Linear algebra fundamentals
    \item Financial markets knowledge
\end{itemize}

% ==================== MODULE STRUCTURE ====================
\section{Module Structure}

\weekheader{Week 1: Foundations of Machine Learning}

\textbf{Topics:} Supervised vs.\ unsupervised learning, loss functions and empirical risk minimization, bias-variance decomposition, model selection (training/validation/test sets), cross-validation for time-series, multiple testing and data snooping in finance.

\textbf{Readings:}
\begin{itemize}[nosep]
    \recommended{James et al.\ (2023) \textit{ISL} -- Chapters 2, 5}
    \recommended{Harvey, Liu \& Zhu (2016) ``...and the Cross-Section of Expected Returns'' \textit{RFS}}
    \advanced{Gu, Kelly \& Xiu (2020) ``Empirical Asset Pricing via ML'' \textit{RFS} -- Intro \& Sec.\ 2}
\end{itemize}

\textbf{Tutorial 1:} Environment setup, data pipeline construction, time-series train-test splits

\weekheader{Weeks 2--3: Regression Methods}

\textbf{Topics:} OLS limitations in high dimensions; Ridge, Lasso, and Elastic Net regularization; tree-based methods (CART, Random Forests, Gradient Boosting); XGBoost/LightGBM; hyperparameter tuning via cross-validation.

\textbf{Application:} Factor investing and market timing using macroeconomic predictors and firm characteristics.

\textbf{Readings:}
\begin{itemize}[nosep]
    \recommended{James et al.\ (2023) \textit{ISL} -- Chapters 3, 6, 8}
    \recommended{Goyal, Welch \& Zafirov (2024) ``A Comprehensive 2022 Look at Equity Premium Prediction'' \textit{RFS}}
    \advanced{Kozak, Nagel \& Santosh (2020) ``Shrinking the Cross-Section'' \textit{JFE}}
    \advanced{Breiman (2001) ``Random Forests'' \textit{Machine Learning}}
\end{itemize}

\textbf{Tutorials 2--3:} Regularized regression pipeline, Random Forest and XGBoost implementation, feature importance analysis

\weekheader{Week 4: Classification Methods}

\textbf{Topics:} Logistic regression; Linear and Quadratic Discriminant Analysis; classification trees and ensembles; evaluation metrics (accuracy, precision, recall, F1, ROC, AUC); class imbalance handling (SMOTE, class weights); probability calibration; threshold optimization.

\textbf{Application:} Credit risk assessment using LendingClub data---default prediction, credit scoring, regulatory considerations.

\textbf{Readings:}
\begin{itemize}[nosep]
    \recommended{James et al.\ (2023) \textit{ISL} -- Chapters 4, 8}
    \recommended{Shumway (2001) ``Forecasting Bankruptcy More Accurately'' \textit{Journal of Business}}
    \recommended{Khandani, Kim \& Lo (2010) ``Consumer Credit-Risk Models via ML'' \textit{JBF}}
    \advanced{Campbell, Hilscher \& Szilagyi (2008) ``In Search of Distress Risk'' \textit{JF}}
\end{itemize}

\textbf{Tutorial 4:} Logistic regression and XGBoost for credit scoring, class imbalance handling, ROC analysis

\weekheader{Weeks 5--7: Advanced Topics}

\textbf{Week 5: Unsupervised Learning}
\begin{itemize}[nosep]
    \item Principal Component Analysis, factor extraction, Instrumented PCA
    \item Clustering methods (K-means, hierarchical)
    \item Applications: dimensionality reduction, portfolio construction, regime detection
\end{itemize}

\textbf{Week 6: Model Interpretability}
\begin{itemize}[nosep]
    \item SHAP values, LIME, partial dependence plots
    \item Walk-forward validation and backtesting
    \item Model monitoring and regulatory compliance (SR 11-7)
\end{itemize}

\textbf{Week 7: Neural Networks}
\begin{itemize}[nosep]
    \item Multi-layer perceptron architecture, activation functions
    \item Backpropagation, regularization (dropout, early stopping)
    \item Applications: return prediction and credit risk comparison
\end{itemize}

\textbf{Readings:}
\begin{itemize}[nosep]
    \recommended{James et al.\ (2023) \textit{ISL} -- Chapters 10, 12}
    \recommended{Kelly, Pruitt \& Su (2019) ``Characteristics are Covariances'' \textit{JFE}}
    \recommended{Lundberg \& Lee (2017) ``A Unified Approach to Interpreting Model Predictions'' \textit{NeurIPS}}
    \advanced{Goodfellow, Bengio \& Courville (2016) \textit{Deep Learning} -- Chapters 6--7}
\end{itemize}

%\textbf{Tutorials 2--8:} PCA and clustering, SHAP analysis, MLP implementation, comprehensive model comparison
\newpage 

% ==================== ASSESSMENT ====================
\section{Assessment}

\subsection{Participation (10\%)}
\begin{itemize}[leftmargin=*]
    \item Active engagement in lectures and tutorials
    \item Quality of questions and contributions
\end{itemize}

\subsection{Group Project (25\%)}

\textbf{Teams of 4--5 students.} Choose one option:

\begin{itemize}[leftmargin=*]
    \item \textbf{Option A: ML Factor Investing Strategy}
    \begin{itemize}
        \item Implement 3+ ML methods for return prediction
        \item Compare to benchmark models
        \item Portfolio construction and backtesting
    \end{itemize}
    
    \item \textbf{Option B: Credit Risk Model}
    \begin{itemize}
        \item Build classification pipeline with 3+ methods
        \item Address class imbalance rigorously
        \item Model interpretation and validation
    \end{itemize}
\end{itemize}

\textbf{Deliverables:} Jupyter notebook, written report (5 pages max), presentation (10 min including Q\&A).

\textbf{Grading:} Implementation quality (40\%), Interpretation and analysis (40\%), Deployment plan (20\%).

\subsection{Final Exam (65\%)}

3-hour written exam with three sections:

\begin{description}[leftmargin=*]
    \item[Section A (30 points):] Conceptual Understanding---when to use which method and why, interpreting model outputs, identifying methodological errors, deployment considerations.
    \item[Section B (20 points):] Applied Problems---interpreting scikit-learn outputs, short coding exercises (fill-in-the-blank), debugging exercises, hyperparameter selection with justification.
    \item[Section C (15 points):] Case Study Analysis---critical evaluation of ML study in finance, identify methodological issues and suggest improvements.
\end{description}

% ==================== TEXTBOOKS ====================
\section{Core Textbooks}

\textbf{Recommended (all available free online):}

\begin{enumerate}[leftmargin=*]
    \item \textbf{James, Witten, Hastie, Tibshirani \& Taylor (2023)} \textit{An Introduction to Statistical Learning with Applications in Python}\\
    Primary textbook, accessible treatment. \url{https://www.statlearning.com/}
    
    \item \textbf{Hastie, Tibshirani \& Friedman (2009)} \textit{Elements of Statistical Learning}\\
    More advanced theoretical reference. \url{https://hastie.su.domains/ElemStatLearn/}
    
    \item \textbf{Goodfellow, Bengio \& Courville (2016)} \textit{Deep Learning}\\
    Neural networks reference. \url{https://www.deeplearningbook.org/}
\end{enumerate}

\textbf{Optional (purchase):}
\begin{itemize}[leftmargin=*]
    \item \textbf{L\'{o}pez de Prado (2018)} \textit{Advances in Financial Machine Learning}---Practical financial ML focus.
\end{itemize}

\subsection{Reading Strategy}

Each lecture has a structured reading list:
\begin{itemize}[leftmargin=*]
    \item \textbf{Recommended readings ($\sim$2 per lecture):} Complement lecture slides and notes; discussed in class; important background.
    \item \textbf{Advanced readings (1--2 per lecture):} For research-interested students; technical proofs and extensions; reference for assignments.
\end{itemize}

Readings are organized thematically, not strictly by week.

% ==================== SOFTWARE ====================
\section{Software and Tools}

\subsection{Required}
\begin{itemize}[leftmargin=*]
    \item Python 3.8+
    \item Core libraries: numpy, pandas, scikit-learn, matplotlib, seaborn
    \item ML libraries: XGBoost, SHAP
    \item IDE: VS Code, Spyder, or PyCharm
\end{itemize}

\subsection{Recommended}
\begin{itemize}[leftmargin=*]
    \item Environment: Anaconda/Miniconda
    \item Version control (optional): GitHub/GitLab
\end{itemize}

\subsection{Datasets}
\begin{itemize}[leftmargin=*]
    \item Return prediction: Goyal and Welch (provided/provided)
    \item Credit risk: LendingClub loan data (public/provided)
\end{itemize}

% ==================== POLICIES ====================
\section{Course Policies}

\subsection{Academic Integrity}
\begin{itemize}[leftmargin=*]
    \item All individual work must be your own
    \item \textbf{AI tools:} May be used for learning, but you must understand all the material. Domain expertise is key to boosting productivity via AI tools. Exam tests independent understanding.
    \item \textbf{Group project:} Collaboration within groups only
    \item \textbf{Proper citations} required for all sources
    \item Plagiarism will result in course failure
\end{itemize}

\subsection{Collaboration Policy}
\begin{itemize}[leftmargin=*]
    \item Tutorials: Collaboration encouraged
    \item Group project: Within groups only
    \item Exam: Individual work only
\end{itemize}

\subsection{Late Submissions}
\begin{itemize}[leftmargin=*]
    \item Group projects: 10\% penalty per day late
    \item Extensions: Only for documented emergencies---contact the Programme Office as soon as possible
    \item All team members share responsibility for timely submission
\end{itemize}

% ==================== DISTINCTIVE FEATURES ====================
\section{What Makes This Module Different}

\begin{enumerate}[leftmargin=*]
    \item \textbf{Production focus:} Not just building models, but deploying them---model monitoring and drift detection, documentation and versioning, deployment considerations throughout.
    
    \item \textbf{Financial context:} Finance-specific challenges---time-series validation (avoid look-ahead bias), overfitting dangers in finance, regulatory requirements.
    
    \item \textbf{Interpretability emphasis:} Black-box models are not enough---SHAP values, LIME, partial dependence plots; explaining predictions to stakeholders.
    
    \item \textbf{AI-assisted learning:} Reflecting modern development practices---learn to use AI tools effectively, but maintain deep understanding.
\end{enumerate}

% ==================== ADDITIONAL RESOURCES ====================
% \section{Additional Resources}

% \subsection{Online Courses}
% \begin{itemize}[leftmargin=*]
%     \item Andrew Ng's Machine Learning (Coursera)
%     \item Fast.ai Practical Deep Learning
%     \item Kaggle Learn tutorials
% \end{itemize}

% \subsection{Practice Platforms}
% \begin{itemize}[leftmargin=*]
%     \item Kaggle competitions
%     \item QuantConnect for algorithmic trading
%     \item Google Colab for free GPU access
% \end{itemize}

% \subsection{Communities}
% \begin{itemize}[leftmargin=*]
%     \item Cross Validated (Stack Exchange)
%     \item Quantitative Finance Stack Exchange
% \end{itemize}

\vfill
\begin{center}
\textit{Last Updated: \today}\\
\textit{Version 2.0}
\end{center}

\end{document}